\documentclass{article}
\usepackage{amsmath,amssymb}
% Copyright (c) 2026 Alberto Villa Osorno.
% SPDX-License-Identifier: MIT
% Shared notation only; this file is included by buildable math documents.

\newcommand{\TritSet}{\{0,1,2\}}
\newcommand{\WordModulus}[1]{3^{#1}}
\newcommand{\WordDomain}[1]{\mathcal{W}_{#1}}
\newcommand{\MemoryState}[1]{\mathcal{M}_{#1}}
\newcommand{\CrazyDigit}{\gamma}
\newcommand{\CrazyN}[3]{\Gamma_{#1}\!\left(#2,#3\right)}
\newcommand{\RotateN}[2]{\rho_{#1}\!\left(#2\right)}
\newcommand{\DecodeOp}{\delta}
\newcommand{\EncryptOp}{\varepsilon}
\newcommand{\LoadOp}{\mathsf{Load}}
\newcommand{\LowerOp}{\mathsf{Lower}}
\newcommand{\VerifyOp}{\mathsf{Verify}}
\newcommand{\ObservedEquivalent}[1]{\equiv_{\mathrm{obs},#1}}
\newcommand{\CostVector}[1]{\mathbf{K}\!\left(#1\right)}
\newcommand{\Transition}[1]{\xrightarrow{#1}}

\title{Search pruning and state canonicalization Research Model}
\begin{document}
\maketitle
Let $C$ be the candidate space, $V(c,p)\in\{0,1\}$ the independent verifier
for candidate $c$ under profile $p$, and $\mathbf{K}(c)$ the measured cost
vector. Search may rank or generate candidates arbitrarily, but only
\[
  C_V = \{c\in C\mid V(c,p)=1\}
\]
is eligible for comparison. Optimization therefore selects verified candidates
under explicit resource budgets rather than treating the search heuristic as a
correctness authority.

For the first conservative pre-identity pruning rule, let $b_i$ be the complete
byte sequence of raw candidate $i$. Define
\[
  i \sim_{\mathrm{bytes}} j \iff b_i=b_j.
\]
The canonical representative of index $i$ is the least earlier index in its exact
byte-equivalence class,
\[
  r(i)=\min\{j\mid j\le i\land b_j=b_i\}.
\]
The retained set is $R=\{i\mid r(i)=i\}$, so the number of required downstream
candidate evaluations is $|R|\le |C|$. Equality of hashes, prefixes, lengths, or
other reduced descriptors is not sufficient for $\sim_{\mathrm{bytes}}$. Thus
this rule is exact but provides zero reduction when every candidate byte sequence
is unique.

Moreover, this relation is maximal for a hardware-neutral pre-verifier layer.
For any two byte-distinct candidates $b_i\ne b_j$, define a permitted trusted
verifier $V_i$ that accepts exactly payload $b_i$ and rejects $b_j$. Then
\[
  V_i(b_i,p)=1 \qquad\text{and}\qquad V_i(b_j,p)=0.
\]
Therefore any generic relation $\approx$ that identifies some
$b_i\ne b_j$ cannot preserve acceptance for every possible trusted verifier.
A relation coarser than $\sim_{\mathrm{bytes}}$ is sound only after adding
profile-, semantics-, or verifier-specific proof obligations to the trusted
boundary.
\end{document}
